% answers_short.tex —— 2005 年·春季北京卷·理科 · 简短答案速查
% 由 tools/make_short_answers.py 自 answers.tex 的 \gksolhead/\gkanswer 提取，
% 并与 2005/2005_数学真题_春季北京理_答案要点.md 逐题比对后生成，勿手改（改源 MD 后重跑）。

\begin{gknotice}
\textbf{使用说明}：本节只给最终答案便于快速核对，完整推导见解析版。答案为编写者独立推导，\textbf{非官方标准答案}。
\end{gknotice}

\gkpart{选择题}
\begin{center}\begin{tabular}{ccccccccc}
\toprule
题号 & 1 & 2 & 3 & 4 & 5 & 6 & 7 & 8 \\
\midrule
答案 & D & A & C & A & B & C & B & A \\
\bottomrule
\end{tabular}\end{center}

\gkpart{填空题}
\begin{itemize}[leftmargin=2.6em]
  \item[\textbf{9.}] $\frac12$
  \item[\textbf{10.}] $\sin\theta=\frac13$，$\cos2\theta=\frac79$
  \item[\textbf{11.}] $m=\sqrt3$
  \item[\textbf{12.}] $(4+2\sqrt2)a^2$
  \item[\textbf{13.}] 二次函数 18 个；一次函数 6 个
  \item[\textbf{14.}] 第一空 $(-4,0)$；第二空 $(-\infty,-6]\cup[2,+\infty)$
\end{itemize}

\gkpart{解答题}
\begin{itemize}[leftmargin=2.6em]
  \item[\textbf{15.}] (1) $M=(3/2,+\infty)$，$N=(-\infty,1)\cup[3,+\infty)$；(2) $M\cap N=[3,+\infty)$，$M\cup N=(-\infty,1)\cup(3/2,+\infty)$
  \item[\textbf{16.}] (1) $\pi/3$；(2) $\arccos\frac58$
  \item[\textbf{17.}] (1) $b_n=3n-1$；(2) $S_n=\frac{n(3n+1)}2$；(3) 当 $1\le n\le18$ 时 $P_n<Q_n$；$n=19$ 时 $P_n=Q_n$；$n\ge20$ 时 $P_n>Q_n$
  \item[\textbf{18.}] (1) $\frac xa+\frac yb=1$；(2) $\frac1{y_1}+\frac1{y_2}=\frac1b$；(3) 角 $MON=\frac\pi2$
  \item[\textbf{19.}] (1) $v=40$ 千米/小时，最大车流量约为 $11.1$ 千辆/小时；(2) $25<v<64$ 千米/小时
  \item[\textbf{20.}] (1) $f(0)=0$，$f(1)=1$；(2) $k_1<k_2<k_3<k_4<k_5$；(3) $f(x)<x$（$0<x<1$）；(4) $\dfrac{a_4+2a_5-a_2-2a_1}{5(a_1+a_2+a_3+a_4+a_5)}$
\end{itemize}
